\documentclass[12pt,a4paper]{article}

\usepackage[T2A]{fontenc}
\usepackage[utf8]{inputenc}
\usepackage[english,russian]{babel}
\usepackage{amsmath,amssymb}
\usepackage{geometry}
\usepackage{booktabs}
\usepackage{array}
\usepackage{hyperref}

\geometry{left=30mm,right=20mm,top=20mm,bottom=20mm}
\hypersetup{hidelinks}
\setlength{\parindent}{1.25cm}
\setlength{\parskip}{0pt}
\renewcommand{\baselinestretch}{1.2}

\newcommand{\RepositoryURL}{}
\newcommand{\Mid}{\operatorname{mid}}
\newcommand{\Rad}{\operatorname{rad}}
\newcommand{\rank}{\operatorname{rank}}
\newcommand{\sgn}{\operatorname{sgn}}

\begin{document}

\begin{titlepage}
\setlength{\baselineskip}{15pt}
\begin{center}
Санкт-Петербургский политехнический университет Петра Великого\\
Физико-механический институт\\
Высшая школа прикладной математики и вычислительной физики

\vspace{30mm}

Направление подготовки\\
"01.03.02. Прикладная математика и информатика"

\vspace{30mm}

{\fontsize{20}{24}\selectfont Дисциплина "Интервальный анализ"}\\[5mm]
Отчет по лабораторной работе №1\\[2mm]
"По интервальному анализу"
\end{center}

\vspace{43mm}

\begin{flushright}
\begin{minipage}{38mm}
\raggedright
\textbf{Работу\\выполнил:}\\
Медведьков Р. А.\\
Группа:\\
5030102/30201
\end{minipage}
\end{flushright}

\vfill

\begin{center}
Санкт-Петербург
\end{center}
\end{titlepage}

\section{Постановка задачи}

Коэффициенты рассматриваемых систем $Ax=b$ известны с интервальной
неопределённостью. Требуется установить, насколько могут измениться
элементы матрицы, прежде чем её столбцы смогут стать линейно зависимыми.
Необходимо определить наименьшее допустимое значение радиуса и подтвердить достижимость
этой границы конкретной точечной матрицей.

Числовые данные заданы следующими матрицами центров:
\[
C_1=
\begin{pmatrix}
0{,}95&1{,}00\\
1{,}05&1{,}00\\
1{,}10&1{,}00
\end{pmatrix},
\qquad
C_2=
\begin{pmatrix}
1{,}10&0{,}90&1{,}10\\
1{,}40&1{,}00&0{,}80\\
0{,}80&1{,}40&1{,}20
\end{pmatrix}.
\]
Для прямоугольного случая предусмотрены два шаблона радиусов:
\[
\delta
\begin{pmatrix}1&1\\1&1\\1&1\end{pmatrix}
\quad\text{и}\quad
\delta
\begin{pmatrix}1&0\\1&0\\1&0\end{pmatrix}.
\]
В квадратном случае радиус каждого элемента $C_2$ составляет $\delta$.
Параметр неотрицателен. Помимо нижней границы требуется указать все
значения $\delta$, при которых в заданном множестве матриц встречается
реализация без полного столбцового ранга.

Упоминание определителя в прямоугольной части задания нуждается
в уточнении: для размера $3\times2$ он не определён. Поэтому здесь
проверяется условие $\rank A'<2$, равносильное одновременному
обнулению всех трёх миноров второго порядка одной матрицы.

\section{Теоретическая часть}

\subsection{Интервальная матрица и критерий особенности}

Разделим описание неопределённости на матрицу центров
$C=\Mid\boldsymbol A$ и радиусы $\Rad\boldsymbol A=\delta D$.
Неотрицательная матрица $D$ определяет, какие коэффициенты могут
изменяться. Множество всех допустимых реализаций имеет вид
\begin{equation}
\boldsymbol A(\delta)=
\bigl\{B=C+E:\ |E_{ij}|\leq\delta D_{ij}\ \text{для всех }i,j\bigr\}.
\label{eq:interval-matrix}
\end{equation}
Выбор поправки одного элемента не ограничивает выбор остальных.
Особенность интервальной матрицы означает, что среди её реализаций
найдётся хотя бы одна вырожденная. Для квадратной реализации это
выражается равенством $\det B=0$, для рассматриваемой прямоугольной
--- неравенством $\rank B<2$.

Вместо поиска по всем элементам $B$ можно искать ненулевой вектор,
который допустимое возмущение переводит в ядро. Согласно критерию
из~\cite[раздел 1.2.5]{book}, это приводит к неравенству
$|Cx|\leq\delta D|x|$. Его применимость к прямоугольному случаю
следует из того же рассуждения:
\begin{equation}
\exists B\in\boldsymbol A(\delta):\ \rank B<n
\quad\Longleftrightarrow\quad
\exists x\in\mathbb R^n\setminus\{0\}:\ 
|Cx|\leq\delta D|x|.
\label{eq:criterion}
\end{equation}
Все сравнения в этом критерии выполняются по компонентам.
Необходимость видна из равенства $(C+E)x=0$:
\[
|(Cx)_i|=|(Ex)_i|\leq\sum_j\delta D_{ij}|x_j|.
\]
Правая часть задаёт максимальную величину, на которую разрешённые
поправки способны изменить результат умножения в строке $i$.
Достаточность можно доказать явным выбором этих поправок.

Обозначим доступную величину изменения через
$q_i=\sum_j\delta D_{ij}|x_j|$. В строках с $q_i>0$ положим
\begin{equation}
E_{ij}=
-\frac{(Cx)_i}{q_i}\,\delta D_{ij}\sgn(x_j).
\label{eq:construction}
\end{equation}
Критерий гарантирует $|(Cx)_i/q_i|\leq1$, поэтому выбранные
элементы $E$ укладываются в разрешённые радиусы. При умножении
на $x$ сумма поправок в строке равна $-(Cx)_i$.
Если $q_i=0$, невязка $(Cx)_i$ уже должна быть нулевой;
такую строку изменять не нужно. Следовательно, $C+E$ принадлежит
множеству~\eqref{eq:interval-matrix} и имеет ненулевое ядро.

\subsection{Почему прямое вычисление интервального определителя недостаточно}

Для квадратной реализации нулевой определитель является точным
признаком вырождения. Но вычисление формулы определителя над
интервалами и нахождение его истинного диапазона --- разные задачи.
Обычная интервальная подстановка даёт гарантированное включение,
которое может оказаться шире множества достижимых значений.

Причина видна на слагаемых $b_{11}b_{22}b_{33}$ и
$-b_{11}b_{23}b_{32}$. Элемент $b_{11}$ в одной реализации имеет
одно значение, тогда как отдельное оценивание слагаемых допускает
разные значения при его повторных вхождениях. Возникает проблема
зависимости интервальных величин.
Поэтому нуль внутри полученного интервала может быть следствием
избыточной ширины оценки, а не реально достижимым определителем.

Для исходной матрицы $C_2$ это можно проверить при $\delta=1/20$.
Интервальное суммирование шести произведений даёт
$[-2507/4000,\ 5931/4000]$, включая нуль. При этом точный диапазон
составляет $[191/1000,\ 711/1000]$: все реализации невырождены.
Точные границы находятся по $2^9$ угловым матрицам, поскольку
при фиксации остальных элементов определитель аффинен по каждому
оставшемуся элементу. Таким образом, использование только
интервального расширения способно занизить искомый порог.
Критерий~\eqref{eq:criterion} устраняет эту неоднозначность:
он требует существования общего ненулевого вектора ядра
для одной допустимой матрицы.

\subsection{Случай одинаковых радиусов всех элементов}

При единичном шаблоне радиусов произведение $D|x|$ состоит
из одинаковых компонент $\sum_j|x_j|=\|x\|_1$.
Критерий однороден по $x$, поэтому масштаб вектора на результат
не влияет. Выберем нормировку $\|x\|_1=1$ и получим задачу
\begin{equation}
\delta_*=\min_{\|x\|_1=1}\|Cx\|_\infty.
\label{eq:uniform}
\end{equation}
Существование решения следует из непрерывности нормы на компактном
множестве нормированных векторов.

Модули можно раскрыть после выбора знаковой области.
Для этого представим компоненты как $x_j=s_jy_j$, где
$s_j\in\{-1,1\}$, $y_j\geq0$. При фиксированных $s_j$
поиск сводится к линейной задаче:
\begin{equation}
\begin{aligned}
&\text{минимизировать } \delta,\\
&y_j\geq0,\qquad \sum_{j=1}^{n}y_j=1,\qquad \delta\geq0,\\
&-\delta\leq\sum_{j=1}^{n}C_{ij}s_jy_j\leq\delta,
\qquad i=1,\ldots,m.
\end{aligned}
\label{eq:lp}
\end{equation}
Для каждой из $2^n$ знаковых областей достаточно рассмотреть
вершины допустимого многогранника: среди них есть точка,
в которой линейная целевая функция достигает минимума.
Сопоставление результатов всех областей даёт глобальное решение
задачи~\eqref{eq:uniform}.

Если $C$ квадратная и обратима, замена $v=Cx$ позволяет получить
дополнительную формулу для проверки порога:
\begin{equation}
\delta_*=
\frac{1}{\|C^{-1}\|_{\infty\to1}},
\qquad
\|C^{-1}\|_{\infty\to1}
=\max_{\|v\|_\infty=1}\|C^{-1}v\|_1.
\label{eq:inverse-norm}
\end{equation}
Норма возмущения здесь согласована с поэлементными радиусами:
$\|E\|_{1\to\infty}=\max_{i,j}|E_{ij}|$.
Применять~\eqref{eq:inverse-norm} в таком виде можно к обратимой
квадратной матрице при одинаковых независимых радиусах всех элементов.

\subsection{Случай одного изменяемого столбца}

При неизменном единичном втором столбце линейная зависимость
означает, что первый столбец постоянен по строкам.
Обозначив его общее значение через $t$, а исходные наблюдения
через $c_i=(C_1)_{i1}$, получаем ограничения $|c_i-t|\leq\delta$.
Иными словами, $t$ должно принадлежать сразу всем трём интервалам
$[c_i-\delta,c_i+\delta]$. Наименьший радиус определяется
половиной разброса наблюдений:
\begin{equation}
\delta_*^{\mathrm{col}}=
\frac{\max_i c_i-\min_i c_i}{2},
\qquad
t_*=\frac{\max_i c_i+\min_i c_i}{2}.
\label{eq:single-column}
\end{equation}
В программной реализации та же граница находится перебором вершин
линейной задачи с переменными $(t,\delta)$, ограничениями
$-\delta\leq c_i-t\leq\delta$ и условием $\delta\geq0$.

\section{Реализация}

Расчёты выполняет файл \texttt{lab1.py}, для запуска достаточно стандартной библиотеки.
Исходные десятичные значения преобразуются в рациональные числа
\texttt{fractions.Fraction}, а комбинации ограничений и знаков
перебираются средствами \texttt{itertools}. Поэтому проверка
равенств проводится точно, без выбора численного допуска.

Сначала для очередного набора знаков составляется система
\eqref{eq:lp}. Затем выбирается столько активных неравенств,
сколько нужно вместе с нормировкой для определения вершины.
Выбранные ограничения рассматриваются как равенства; решение
полученной системы находится исключением Гаусса. Вырожденные
системы пропускаются. Остальные кандидаты проверяются по полному
набору ограничений, после чего совпадающие вершины внутри одной
знаковой области удаляются. Из общего списка допустимых точек
выбирается решение с наименьшим радиусом.

Для фиксированного второго столбца используется тот же принцип,
но вершины ищутся непосредственно в плоскости $(t,\delta)$.
Найденный вектор ядра позволяет восстановить точечную матрицу
по~\eqref{eq:construction}. Далее программа отдельно проверяет
её элементы, нулевую невязку, ранг и максимальные миноры.
Контроль квадратного случая выполняется также через
\eqref{eq:inverse-norm}; пример с радиусом $1/20$ проверяется
сравнением интервального расширения и точного диапазона определителя.

Команда запуска: \texttt{python lab1.py}. Русскоязычный протокол
сохраняется в \texttt{results/console.txt}, а дроби и результаты
проверок --- в \texttt{results/results.json}.
Файл отчёта самодостаточен для компиляции в pdfLaTeX.

Выбранный перебор удобен для малых размерностей $3\times2$ и
$3\times3$: он даёт точный результат и позволяет проверить каждый
кандидат. При увеличении числа переменных вычислительные затраты
быстро растут, поэтому данный способ не рассчитан на большие матрицы.

\section{Результаты}

\subsection{Прямоугольная матрица: изменяются оба столбца}

Минимум, найденный перебором для $C_1$, составляет
$\delta_*=1/27\approx0{,}037037$. Убедимся аналитически,
что меньший радиус невозможен. Если $x_1\ne0$, введение
$t=-x_2/x_1$ преобразует критерий в оценку
\[
\delta\geq
\frac{\max\{|0{,}95-t|,\ |1{,}10-t|\}}{1+|t|}.
\]
Отрицательные $t$ дают правую часть больше единицы и потому
не улучшают найденный результат. На неотрицательной полуоси
минимум достигается при $t=41/40$: до этой точки отношение
убывает, а после неё возрастает. Подстановка даёт
\[
\frac{3/40}{1+41/40}=\frac1{27}.
\]
Оставшийся случай $x_1=0$ возможен по критерию только при
$\delta\geq1$. Таким образом, все направления $x$ учтены.

Вектор, выбранный программой на границе, равен
$x=(-40/81,\ 41/81)^\mathsf T$. Ему соответствует реализация
\[
A'_{1,\mathrm{all}}=
\begin{pmatrix}
\frac{533}{540}&\frac{26}{27}\\[2pt]
\frac{1681}{1620}&\frac{82}{81}\\[2pt]
\frac{287}{270}&\frac{28}{27}
\end{pmatrix}.
\]
Наибольшее отклонение её элемента от центра составляет ровно
$1/27$, то есть интервальные ограничения соблюдены.
Столбцы связаны соотношением $A'_{:1}=(41/40)A'_{:2}$.
Отсюда $A'_{1,\mathrm{all}}x=0$ и ранг равен единице;
все три минора второго порядка также обращаются в нуль.

\subsection{Прямоугольная матрица: изменяется первый столбец}

Разброс элементов первого столбца равен $1{,}10-0{,}95=0{,}15$.
Поэтому формула~\eqref{eq:single-column} даёт
$\delta_*^{\mathrm{col}}=0{,}15/2=3/40$.
Первое общее значение трёх интервалов равно $t_*=41/40=1{,}025$.
При этом получаем следующую матрицу и ненулевой вектор её ядра:
\[
A'_{1,\mathrm{col}}=
\begin{pmatrix}
\frac{41}{40}&1\\
\frac{41}{40}&1\\
\frac{41}{40}&1
\end{pmatrix},
\qquad
x=\begin{pmatrix}1\\-\frac{41}{40}\end{pmatrix}.
\]
Чтобы получить эту реализацию, к первому столбцу прибавляются
поправки $(3/40,-1/40,-3/40)^\mathsf T$.
Они допустимы при найденном радиусе, а элементы второго столбца
сохранены точно. Первый столбец есть второй, умноженный на $41/40$,
так что $\rank A'_{1,\mathrm{col}}=1$ и $A'_{1,\mathrm{col}}x=0$.
Это подтверждает достижимость нижней границы внутри
$\boldsymbol A_1(3/40)$.

\subsection{Квадратная матрица}

Перебор опорных решений квадратной задачи даёт порог $1/10$.
Для его независимого подтверждения вычислим норму из
\eqref{eq:inverse-norm}. Обратная матрица в точной арифметике равна
\[
C_2^{-1}=\frac1{107}
\begin{pmatrix}
20&115&-95\\
-260&110&165\\
290&-205&-40
\end{pmatrix}.
\]
Функция $v\mapsto\|C_2^{-1}v\|_1$ выпукла, поэтому её максимум
на кубе $\|v\|_\infty\leq1$ достигается среди вершин.
Их восемь, но противоположные вершины имеют одинаковые значения
нормы. Достаточно следующих четырёх вычислений:
\[
\begin{array}{c|cccc}
v&(1,1,1)&(1,1,-1)&(1,-1,1)&(1,-1,-1)\\ \hline
\|C_2^{-1}v\|_1&100/107&670/107&850/107&10
\end{array}
\]
Наибольшее значение равно $10$. Его обратная величина $1/10=0{,}1$
совпадает с результатом основного алгоритма.

Построенная по найденному вектору ядра реализация имеет вид
\[
A'_2=
\begin{pmatrix}
\frac{11}{10}&1&1\\
\frac75&\frac9{10}&\frac9{10}\\
\frac45&\frac{13}{10}&\frac{13}{10}
\end{pmatrix},
\qquad
x=\begin{pmatrix}0\\-\frac12\\\frac12\end{pmatrix}.
\]
Изменения коэффициентов второго и третьего столбцов по модулю
равны $1/10$, а первый столбец остаётся исходным.
Таким образом, $A'_2$ лежит в заданных интервалах.
Равенство последних двух столбцов даёт $A'_2x=0$ и $\det A'_2=0$.
Минор первых двух строк и первых двух столбцов равен $-41/100\ne0$,
поэтому ранг составляет ровно два.

\begin{table}[ht]
\centering
\caption{Границы особенности и характеристики найденных реализаций}
\label{tab:results}
\begin{tabular}{>{\raggedright\arraybackslash}p{49mm}ccc}
\toprule
Вариант & Диапазон $\delta$ & Ранг $A'$ & Опорных решений\\
\midrule
$A_1$, изменяются оба столбца
  & $[1/27,\,+\infty)$ & 1 & 10\\
$A_1$, изменяется первый столбец
  & $[3/40,\,+\infty)$ & 1 & 1\\
$A_2$, изменяются все элементы
  & $[1/10,\,+\infty)$ & 2 & 58\\
\bottomrule
\end{tabular}
\end{table}

В последнем столбце таблицы~\ref{tab:results} приведено суммарное
количество допустимых вершин по всем знаковым областям.
Внутри области повторы удаляются, но точка с нулевыми компонентами
может учитываться в разных областях. Поэтому эти числа не следует
трактовать как количество различных матриц-свидетелей.

\section{Обсуждение}

Найденные границы разделяют два режима. Ниже порога все допустимые
реализации имеют полный столбцовый ранг. На границе впервые
становится возможной линейная зависимость столбцов.
Дальнейшее увеличение радиуса не устраняет эту возможность:
построенная реализация остаётся внутри расширившихся интервалов.
Именно поэтому каждый диапазон в таблице~\ref{tab:results}
начинается с найденного минимума и не ограничен сверху.
При этом невырожденные реализации также сохраняются, в частности
сама матрица центров. Речь идёт о возможности потери ранга,
а не о её обязательности для всех реализаций.

Регрессионная интерпретация связана с моделью $b_i\approx k t_i+c$:
первый столбец содержит аргументы наблюдений, а единичный второй
представляет свободный член. При радиусе $3/40$ аргументы могут
стать одинаковыми и равными $41/40$. Тогда по данным различима
лишь комбинация $(41/40)k+c$, а раздельно определить $k$ и $c$
по такой матрице уже нельзя.

В томографической интерпретации каждая строка описывает измерение,
а каждый столбец --- вклад одной неизвестной компоненты объекта.
Если разрешить независимые изменения обоих столбцов $A_1$,
их пропорциональность достижима при меньшем радиусе $1/27$.
Фиксация единичного столбца повышает порог: для данных чисел
отношение составляет $(3/40)/(1/27)=81/40=2{,}025$.
В квадратном примере на границе совпадают вклады второй и третьей
компонент, так что изменения вдоль $(0,-1,1)^\mathsf T$
не отражаются на измерениях.

Поскольку значения правой части $b$ не заданы, выводы касаются
только ранга матрицы. Потеря ранга сама по себе не гарантирует
совместность системы. Если решение существует, к нему можно
прибавлять вектор ядра; при неподходящей правой части решений
может не оказаться.

\section{Выводы}

Перебор знаковых областей и активных ограничений позволил
определить точные пороги для всех вариантов задания.
При неопределённости обоих столбцов $A_1$ граница равна $1/27$,
при фиксированном втором столбце --- $3/40$, а для $A_2$ --- $1/10$.
Во всех случаях особенные реализации существуют ровно при
$\delta\geq\delta_*$, включая граничное значение.

Достижимость каждого порога подтверждена построением матрицы
с ненулевым ядром. Проверки элементов и максимальных миноров
согласуются с полученными рангами: один для прямоугольных
реализаций и два для квадратной. Рациональная арифметика
обеспечила точное выполнение проверяемых равенств.
Аналитические оценки и норма обратной матрицы независимо
подтвердили минимальность вычисленных радиусов.

\section{Список литературы}

\begingroup
\makeatletter
\renewcommand{\section}[2]{}
\makeatother
\begin{thebibliography}{9}
\bibitem{book}
Баженов А.~Н., Карпова А.~А.
\textit{Интервальный анализ в примерах}: учебное пособие.
Санкт-Петербург, 2024. 226~с.


\end{thebibliography}
\endgroup


\end{document}
